\documentclass[11pt,oneside]{report}
% ======================================================================
%  THE TRIADIC MEASUREMENT SUBSTRATE --- COMPLETE CORPUS (Volumes 1 & 2)
%  Aaron Switzer, 2026
%  Single-file build. Compile: pdflatex x3 (third pass settles the TOC).
% ======================================================================
% ======================================================================
%  TMS Volume 1 -- shared preamble
%  Compiles with pdfLaTeX both locally and on Overleaf (no exotic fonts).
% ======================================================================
\usepackage[T1]{fontenc}
\usepackage[utf8]{inputenc}
\usepackage{amsmath}
\usepackage{amssymb}
\usepackage{mathptmx}            % Times text + math (widely available)
\usepackage[scaled=0.90]{helvet} % sans for headings
\usepackage{courier}

\usepackage[letterpaper,margin=1in]{geometry}
\usepackage{amsthm}
\usepackage{mathtools}
\usepackage{bm}
\usepackage{microtype}
\usepackage{enumitem}
\usepackage{booktabs}
\usepackage{array}
\usepackage{longtable}
\usepackage{caption}
\usepackage{xcolor}
\usepackage{titlesec}
\usepackage{fancyhdr}
\usepackage[hidelinks]{hyperref}
\usepackage{tocbibind}

% ---- spacing -----------------------------------------------------------
\setlength{\parindent}{1.2em}
\setlength{\parskip}{0.35em}
\linespread{1.04}
\setlist{leftmargin=1.6em,topsep=0.25em,itemsep=0.15em,parsep=0pt}

% ---- theorem environments ---------------------------------------------
\newtheorem{theorem}{Theorem}[section]
\newtheorem{lemma}[theorem]{Lemma}
\newtheorem{corollary}[theorem]{Corollary}
\newtheorem{proposition}[theorem]{Proposition}

\theoremstyle{definition}
\newtheorem{definition}[theorem]{Definition}
\newtheorem{axiom}{Axiom}
\newtheorem{claim}[theorem]{Claim}
\newtheorem{principle}[theorem]{Principle}
\newtheorem{conjecture}[theorem]{Conjecture}
\newtheorem{prediction}[theorem]{Prediction}
\newtheorem{property}[theorem]{Property}
\newtheorem{postulate}[theorem]{Postulate}

\theoremstyle{remark}
\newtheorem{remark}[theorem]{Remark}
\newtheorem*{remark*}{Remark}

% ---- section heading look (unnumbered; author supplies the numbers) ----
\titleformat{\section}[hang]{\normalfont\large\bfseries\sffamily}{}{0pt}{}
\titleformat{\subsection}[hang]{\normalfont\normalsize\bfseries\sffamily}{}{0pt}{}
\titleformat{\subsubsection}[hang]{\normalfont\normalsize\bfseries\itshape}{}{0pt}{}
\titlespacing*{\section}{0pt}{1.4em}{0.5em}
\titlespacing*{\subsection}{0pt}{1.0em}{0.35em}
\titlespacing*{\subsubsection}{0pt}{0.8em}{0.3em}

% chapter (= each paper) heading
\titleformat{\chapter}[display]
  {\normalfont\sffamily\bfseries}
  {}{0pt}{\Huge}
\titlespacing*{\chapter}{0pt}{10pt}{24pt}

% ---- page-break quality: avoid stranded headings and orphaned/widowed lines ----
% (applies to both volumes via the shared preamble)
\widowpenalty=10000        % no single line at the top of a page
\clubpenalty=10000         % no single line at the bottom of a page
\displaywidowpenalty=10000 % same, around displayed equations
\brokenpenalty=4000        % discourage page breaks right after a hyphenated line
% titlesec: forbid a page break immediately after a heading (heading stranded at page foot)
\usepackage{needspace}
\usepackage{etoolbox}
\pretocmd{\section}{\Needspace*{4\baselineskip}}{}{}
\pretocmd{\subsection}{\Needspace*{3\baselineskip}}{}{}

% ---- running heads -----------------------------------------------------
\pagestyle{fancy}
\fancyhf{}
\fancyhead[L]{\small\itshape The Triadic Measurement Substrate}
\fancyhead[R]{\small\itshape\nouppercase{\rightmark}}
\fancyfoot[C]{\small\thepage}
\renewcommand{\headrulewidth}{0.4pt}
\fancypagestyle{plain}{\fancyhf{}\fancyfoot[C]{\small\thepage}\renewcommand{\headrulewidth}{0pt}}

% ---- abstract / keywords / references list ----------------------------
\newenvironment{paperabstract}
  {\begin{center}\begin{minipage}{0.88\textwidth}\small
   \noindent\textbf{Abstract.}\itshape\space}
  {\end{minipage}\end{center}\vspace{0.4em}}

\newcommand{\keywords}[1]{\noindent{\small\textbf{Keywords:} #1}\vspace{0.4em}}

% per-paper APA reference list (hanging indent), kept as authored
\newenvironment{reflist}
  {\par\small\setlength{\parindent}{0pt}%
   \setlength{\leftskip}{1.6em}\setlength{\parindent}{-1.6em}%
   \setlength{\parskip}{0.4em}%
   \raggedright\hyphenpenalty=50\exhyphenpenalty=50\relax
   \sloppy}
  {\par}

% convenience
\providecommand{\tightlist}{\setlength{\itemsep}{0pt}\setlength{\parskip}{0pt}}
\providecommand{\TMS}{\textsc{tms}}
\hyphenation{con-fir-ma-tion sub-strate}
\begin{document}

\thispagestyle{empty}
\begin{titlepage}
\centering
\vspace*{2cm}
{\sffamily\bfseries\Large The Triadic Measurement Substrate\par}
\vspace{0.6cm}
{\large $\bar{B} = 3\bar{\alpha}$\par}
\vspace{1cm}
{\sffamily\large Volume 1\par}
\vspace{1.4cm}
{\Large \bfseries Paper 2: Computational Dynamics of the Binary Alphabet: The Polarization Field, the Born-Rule Structural Correspondence, and the Emergence of Universal Computation\par}
\vfill
{\large Aaron Switzer\par}
\vspace{0.4cm}
{\large 2026\par}
\vspace{1.2cm}
{\small\itshape Excerpted from the complete two-volume corpus, \emph{The Triadic Measurement Substrate}. Full corpus and reproducibility package: \url{https://doi.org/10.5281/zenodo.22035422}\par}
\end{titlepage}
\cleardoublepage
\pagenumbering{arabic}

% ---------------- Volume 1 / paper2 ----------------
\chapter*{Paper 2}
\markboth{Paper 2}{Paper 2}
\addcontentsline{toc}{chapter}{Paper 2: Computational Dynamics of the Binary Alphabet: The Polarization Field, the Born-Rule Structural Correspondence, and the Emergence of Universal Computation}
\begingroup\centering
{\large\itshape Computational Dynamics of the Binary Alphabet: The Polarization Field, the Born-Rule Structural Correspondence, and the Emergence of Universal Computation\par}\vspace{0.8em}
{\normalsize Aaron Switzer\par}
{\small May 2026\par}
\endgroup\vspace{1.0em}

\noindent{\small\textbf{Keywords:} TMS, Binary Alphabet, Polarization Field, Born Rule, Majority Gate, Universal Computation, Programs, Stabilizer Dynamics, Computational Irreducibility}\par\vspace{0.6em}

\section*{Status of Results}
\addcontentsline{toc}{section}{Status of Results}

This paper distinguishes four categories of claim, each labeled where it
appears in the text.

\textbf{Derived.} Results that follow from the five axioms, the
Principle of Generative Economy, and the structures established in Paper
1 by explicit construction or proof. The polarization conjugate pair \ensuremath{(\psi ,}
\ensuremath{\Pi_{\psi})} and the polarization wave equation \ensuremath{\partial ^{2}\psi /\partial t^{2}} \ensuremath{-} \ensuremath{c^{2}\nabla ^{2}\psi} = 0 (\S\,II), the
propagation speed \ensuremath{c_{\mathrm{TMS}}} = \ensuremath{1/\sqrt{2}} for \ensuremath{\psi} inherited from FCC geometry (\S\,2.3),
the three-way logical distinction among possibilities for synthetic-bit
valuation and the selection of polarization-biased stochastic resolution
(\S\,1.2, \S\,3.3), the linear amplitude-to-probability mapping \ensuremath{P(\bar{B}} = 1
\textbar{} \ensuremath{\Pi_{\psi})} = (1 + \ensuremath{\Pi_{\psi})/2} from four geometric constraints (\S\,4.1),
the 3-input majority gate at every confirmation event (\S\,5.1), and
universal classical computation from \{MAJ, constant\} (\S\,5.3) are
derived in this sense. The PGE-Forced Wave-Form Theorem (\S\,VII):
higher-order reinterpretation at any hierarchical boundary is wave-form
by structural necessity, not by additional axiom. Used in Paper 5 \S\,III
to resolve the two-waves problem and in Paper 5 \S\,IV for the QCA-to-QFT
correspondence. The formal definition of a program as a configuration
causally closed under T with K(C) \textless{} \textbar C\textbar{}
(\S\,6.2) is a definition; the existence of programs is established by
construction in Paper 3.

\textbf{Structural correspondence.} Results in which a TMS construction
faithfully reproduces the skeleton of a known framework without claiming
reproduction of every feature. The Born-rule correspondence of \S\,4.2 is
the central case: the skeleton of ``deterministic wave evolution between
events plus amplitude-biased projective measurement at each event'' is
reproduced exactly; the linear mapping (1 + \ensuremath{\Pi_{\psi})/2} is the unique
geometrically-forced form for two-outcome real-valued projective
measurement. The full quadratic \textbar\ensuremath{\cdot}\textbar\ensuremath{^{2}} mapping requires
complex amplitudes and phase coherence, which are expected to arise from
the macroscopic interference behavior of coupled \ensuremath{(\psi ,} \ensuremath{\Pi_{\psi})} propagation
and are deferred. The label ``structural correspondence'' matches the
[[4,2,2]] correspondence of Paper 1: skeleton reproduced, full
formalism deferred.

\textbf{Predictions / Conjectures.} Falsifiable claims that follow from
the framework but whose closed forms are not derived within this paper.
The refinement \ensuremath{\beta} = f(d, \ensuremath{K_{\psi})} where \ensuremath{K_{\psi}} is the Kolmogorov complexity of
the program running through the bit's causal history (\S\,6.4) is a
sharpening of the Paper 1 conjecture, not a derivation of f.~The
empirical signatures of polarization refraction (\S\,7.1) and the
prediction that GUP corrections vary with a particle's preparation
history rather than being a fixed property of a species (\S\,7.2) are
consequences of the framework's architecture; the explicit functional
forms and falsification thresholds are deferred to Paper 3's refinements
and Paper 5's simulation work.

\textbf{Deferred derivations carried into Paper 3.} Three results are
stated in this paper and closed by derivation in Paper 3: the
dissolution rate of non-T-closed configurations (\S\,6.3, closed in Paper 3
\S\,2.5 as \ensuremath{\sqrt{2}} \ensuremath{\cdot} \ensuremath{L_{R}/\ell_{p}} flops); the effective error-rate decay with
coherence depth (\S\,5.2, closed in Paper 3 \S\,2.5 as \ensuremath{(1/3)^{d}} in the
large-coherence limit); and the explicit conditions under which the
stabilizer dynamics select program-bearing configurations against \ensuremath{U_{\mathrm{sh}}}
disruption (\S\,6.3, closed in Paper 3 \S\,3.3 as the coherence bistability
threshold). These are flagged here as deferrals rather than completed
derivations.

This four-way labeling applies throughout the paper. Every claim in the
body text falls into exactly one category.

\section*{Preface to Paper 2}
\addcontentsline{toc}{section}{Preface to Paper 2}

This paper is the second installment in Volume 1 of the Triadic
Measurement Substrate program. Volume 1, Paper 1 established the
geometric and ontological architecture of the TMS: the agent manifold
\ensuremath{M\alpha ,} the triadic confirmation rule, the \ensuremath{1\to 4} multiplier, FCC necessity,
the identification of the agent spacing with the Planck length, the
conjugate field structure \ensuremath{(\varphi ,} \ensuremath{\Pi ),} the [[4,2,2]] structural
correspondence, the binary alphabet \ensuremath{\bar{B}} \ensuremath{\in} \{0, 1\} as the irreducible
computational content of every confirmed bit, and the polarization field
\ensuremath{\psi} as a parallel object propagating across the same FCC lattice. Paper 2
takes that architecture as given and develops its computational
dynamics. No new axioms are introduced. No new free parameters are
introduced. The Principle of Generative Economy continues to govern
every selection.

Paper 2's central claim is that the same primitives that generate wave
dynamics in Paper 1 --- by geometric necessity --- also generate biased
stochastic measurement, universal computation, and stabilizer-protected
programs. The Born rule, the majority gate, and computational
irreducibility all emerge from the interaction of the polarization field
\ensuremath{\psi} with the triadic confirmation rule and the maximum-entropy substrate
\ensuremath{U_{\mathrm{sh}}.} Computation in the TMS is not designed. It is the inevitable
consequence of confirming a binary outcome against a
polarization-modulated wave field.

\begin{paperabstract}
Building on the binary alphabet \ensuremath{\bar{B}} \ensuremath{\in} \{0, 1\} and the parallel
polarization field \ensuremath{\psi} established in Volume 1, Paper 1, we develop the
computational dynamics of the TMS substrate. Three principal results are
derived. First, the polarization field \ensuremath{\psi} obeys a homogeneous wave
equation on the FCC lattice identical in form and propagation speed to
that of the confirmation density \ensuremath{\varphi ,} establishing \ensuremath{\psi} as a fully-fledged
conjugate-pair dynamical object propagating at \ensuremath{c_{\mathrm{TMS}}} = \ensuremath{1/\sqrt{2}} lattice
spacings per flop. Second, at every confirmation event the resolution of
a synthetic bit position into a binary value \ensuremath{\bar{B}} \ensuremath{\in} \{0, 1\} is biased by
the local incoming polarization wave amplitude against the unbiased
stochastic kick from \ensuremath{U_{\mathrm{sh}}} --- yielding a linear amplitude-to-probability
mapping \ensuremath{P(\bar{B}} = 1 \textbar{} \ensuremath{\Pi_{\psi})} = (1 + \ensuremath{\Pi_{\psi})/2} that constitutes a
structural correspondence to the Born rule for projective measurement.
Third, the geometric structure of triadic confirmation makes every
confirmation event a probabilistic 3-input majority gate over the
polarizations of its parent confirmations; the [[4,2,2]]
stabilizer dynamics from Paper 1 correct local errors and project the
effective behavior toward deterministic majority logic at scale.
Universal computation follows. We define a program as a configuration of
the polarization landscape that reproduces itself under the flop
operator T while remaining algorithmically compressible (K(C)
\textless{} \textbar C\textbar), and show that the stabilizer dynamics
provide an automatic selection mechanism: configurations that satisfy
stabilizer parity and reproduce under T persist; configurations that do
not dissolve into \ensuremath{U_{\mathrm{sh}}-\text{like}} noise. The \ensuremath{\beta (d,} \ensuremath{K_{\psi})} prediction sharpened in
Paper 1 acquires its computational meaning here: the depth and
polarization-complexity of a bit's causal history is the depth and
complexity of the program running through it. As in Paper 1, no free
numerical parameters are introduced; every result follows by geometric
and informational necessity from the axioms.

Paper 2 engages directly with two active programs deriving quantum
mechanics from substrate dynamics: 't Hooft's Cellular Automaton
Interpretation, which derives apparent quantum stochasticity from a
deterministic underlying automaton observed in the wrong basis ('t
Hooft, 2016), and Gorard's derivation of quantum measurement from causal
invariance and multiway-graph collapse in the Wolfram model (Gorard,
2020). The distinctive move of TMS relative to these programs is the
polarization-biased resolution rule on a single confirmed lattice, which
derives the linear amplitude-to-probability mapping from signal
conservation and the maximum-entropy character of \ensuremath{U_{\mathrm{sh}}} without invoking
either superposition over multiway branches or a hidden deterministic
preferred basis.
\end{paperabstract}

\section*{I. Foundations from Paper 1 and the Computational Program}
\addcontentsline{toc}{section}{I. Foundations from Paper 1 and the Computational Program}

\subsection*{1.1 Inherited Structure}

Paper 1 establishes the following structures, which Paper 2 takes as
given without re-derivation:

\begin{itemize}
\item
  The agent manifold \ensuremath{M\alpha} with agent spacing \ensuremath{\ell_{0}} \ensuremath{\equiv} \ensuremath{\ell_{p}} (Axiom 2; Section 3.2
  of Paper 1).
\item
  The FCC causal stack with face-diagonal nearest-neighbor connectivity
  (Section 2.4 of Paper 1).
\item
  The triadic actualization rule \ensuremath{\bar{B}} = \ensuremath{3\bar{\alpha}_{2}} in 2D and \ensuremath{\bar{B}} = \ensuremath{4\bar{\alpha}_{3}} in 3D
  (Sections 2.1, 2.5 of Paper 1).
\item
  The \ensuremath{1\to 4} multiplier with informational closure (Section 4 of Paper 1).
\item
  The binary alphabet \ensuremath{\bar{B}} \ensuremath{\in} \{0, 1\} with polarization \ensuremath{\chi} \ensuremath{\equiv} \ensuremath{2\bar{B}} \ensuremath{-} 1 \ensuremath{\in} \{\ensuremath{-1,}
  +1\} (Section 1.4 of Paper 1).
\item
  Polarization-carrying touch vector signals (Section 1.3 of Paper 1).
\item
  The flop as the atomic unit of state change with no sub-flop time
  (Section 1.5 of Paper 1).
\item
  The conjugate field structure \ensuremath{(\varphi ,} \ensuremath{\Pi )} and the wave equation for \ensuremath{\varphi} at
  \ensuremath{c_{\mathrm{TMS}}} = \ensuremath{1/\sqrt{2}} (Section 6 of Paper 1).
\item
  The [[4,2,2]] structural correspondence with logical qubit LQ1
  = bit value, LQ2 = causal-history inheritance, jointly protected by \ensuremath{S_{x}}
  and \ensuremath{S_{Z}} (Section 5 of Paper 1).
\item
  The polarization field \ensuremath{\psi} as parallel propagating object on the same
  FCC lattice (Section 6.8 of Paper 1).
\item
  \ensuremath{U_{\mathrm{sh}}} as the maximum-entropy substrate from which both latent bits and
  stochastic disruptions arise, with no preferred direction, scale, or
  target (Axiom 0; Section 6.6 of Paper 1).
\end{itemize}

These structures are the entirety of the foundation. Paper 2 introduces
no further primitives.

\begin{remark}[Forward-Compatibility with the Sublattice Bifurcation]
At every
hierarchical level k \ensuremath{\geq} 1, the meta-FCC splits into two interpenetrating
sublattices A and B running independent [[4,2,2]] dynamics
(Paper 1 \S\,2.4 corollary, full development in Paper 5 \S\,6.1.3 and \S\,5.3).
This affects no Paper 2 derivations at the substrate level but is
relevant when the polarization-field and Born-rule arguments are
recursively applied at meta-levels.
\end{remark}

\subsection*{1.2 The Computational Question}

Paper 1 closed by establishing that every confirmed bit carries a binary
value. It deferred the question of how that value is determined at the
moment of confirmation. Paper 2 answers it.

The question is sharper than it first appears. The synthetic bit
positions of the next layer \ensuremath{S_{n+1}} are geometrically fixed by the \ensuremath{1\to 4}
multiplier acting on the parent confirmations in \ensuremath{S_{n}.} Every synthetic bit
position is a geometric appointment --- a location where a confirmation
event must occur at the next flop. The question is not whether the
appointment is kept (informational closure forces it to be) but with
what value.

Three logically distinct possibilities exist:

\begin{itemize}
\item
  Inherited value. The synthetic bit takes a deterministic function of
  its parents' values. This requires specifying that function --- a rule
  mapping \ensuremath{(\chi_{1},} \ensuremath{\chi_{2},} \ensuremath{\chi_{3})} \ensuremath{\to} \ensuremath{\chi_{\mathrm{child}}} --- which is an external program
  imposed on the substrate. It violates the Principle of Generative
  Economy.
\item
  Independent randomness. The synthetic bit takes an unbiased random
  value from \ensuremath{U_{\mathrm{sh}},} with no influence from the parents. This violates
  Axiom 3 (signal conservation): the polarizations carried by the 12
  outgoing signals from the parent confirmations are then dissipated
  rather than expressed in the next layer's confirmations.
\item
  Polarization-biased stochastic resolution. The synthetic bit's
  resolution is stochastic, but the probability distribution is biased
  by the local incoming polarization carried by its parents' touch
  vector signals. The bias is geometric. The randomness is the residual
  contribution from \ensuremath{U_{\mathrm{sh}}.}
\end{itemize}

Possibility 3 is the unique selection under the joint constraints of
Axiom 3 and the Principle of Generative Economy. Polarization signals
are conserved by being expressed as bias, not dissipated; no external
rule is imposed; \ensuremath{U_{\mathrm{sh}}} provides the stochastic component as it does at
every scale of the TMS. Sections III and IV formalize this resolution
mechanism and identify it as a structural correspondence to the Born
rule for projective measurement.

\section*{II. The Polarization Field \ensuremath{\psi} and its Wave Equation}
\addcontentsline{toc}{section}{II. The Polarization Field \ensuremath{\psi} and its Wave Equation}

\subsection*{2.1 Definition of the Polarization Field}

\begin{definition}[Polarization Field]
Let \ensuremath{\psi (\bar{r},} \ensuremath{\tau )} denote the local
coarse-grained polarization --- the scalar field representing the mean
polarization of confirmed bits at any point in the FCC lattice:
\end{definition}

\begin{equation*}
\psi (\bar{r}, \tau ) = \langle \chi \rangle (\bar{r}, \tau ) = \langle 2\bar{B} - 1\rangle (\bar{r}, \tau )
\end{equation*}

averaged over a region containing many lattice sites but small compared
to wave scales. \ensuremath{\psi} ranges over \ensuremath{[-1,} +1]. \ensuremath{\psi} = +1 corresponds to a
region of pure confirmed-presence; \ensuremath{\psi} = \ensuremath{-1} corresponds to a region of
pure confirmed-absence; \ensuremath{\psi} = 0 corresponds to a region of balanced binary
content. \ensuremath{\psi} is defined wherever the confirmation density \ensuremath{\varphi} is non-zero;
in regions of vanishing \ensuremath{\varphi ,} \ensuremath{\psi} is undefined.

The polarization field is distinct from the confirmation field. \ensuremath{\varphi} counts
whether interstices are confirmed; \ensuremath{\psi} records the binary content of those
confirmations. The two fields are independent in the bulk: any value of
\ensuremath{\psi} \ensuremath{\in} \ensuremath{[-1,} +1] is compatible with any value of \ensuremath{\varphi} \ensuremath{\in} (0, 1].

\subsection*{2.2 The Conjugate Excitation \ensuremath{\Pi_{\psi}}}

By the same argument as Section 6.3 of Paper 1, the FCC cascade applied
to the polarization field generates a forward-propagating excitation
that must be tracked separately from the field itself. Define:

\begin{equation*}
\Pi_{\psi}(S_{n}) = \text{the} \text{inherited} \text{polarization} \text{pressure} \text{being} \text{driven}
\text{forward} \text{into} S_{n+1}
\end{equation*}

\ensuremath{\psi} records the polarization that has locked in past layers; \ensuremath{\Pi_{\psi}} records
the polarization that the touch vector network is carrying forward to
the next layer. The pair \ensuremath{(\psi ,} \ensuremath{\Pi_{\psi})} is the polarization analogue of the
\ensuremath{(\varphi ,} \ensuremath{\Pi )} pair from Paper 1.

\begin{lemma}[Polarization Conjugacy]
The polarization field \ensuremath{\psi} and the
forward-propagating polarization excitation \ensuremath{\Pi_{\psi}} constitute a
co-generated conjugate pair, arising necessarily from the interaction
between the polarization-carrying \ensuremath{1\to 4} cascade and \ensuremath{U_{\mathrm{sh}}} disruption.
\end{lemma}

\begin{proof}
The argument of Paper 1 Lemma (Field Conjugacy) applies in full,
with three substitutions: (a) confirmation events become polarized
confirmation events; (b) signals become polarized signals; (c) the
conserved quantity is not the count of signals but the signed sum of
their polarizations. Step 1 (2D first-order closure), Step 2 \ensuremath{(1\to 4}
cascade), Step 3 \ensuremath{(U_{\mathrm{sh}}} regulation produces dynamic stability), and Step 4
(signal conservation forces conjugate coupling) carry through unchanged.
By signal conservation applied to polarization content:

\begin{equation*}
\partial \psi / \partial \tau = \Pi_{\psi}
\end{equation*}

\begin{equation*}
\partial \Pi_{\psi} / \partial \tau = c^{2} \Delta_{\mathrm{FCC}} \psi
\end{equation*}

Step 5 (elimination) yields:

\begin{equation*}
\partial ^{2}\psi / \partial \tau^{2} = c^{2} \Delta_{\mathrm{FCC}} \psi
\end{equation*}

The conjugate-pair structure is not an analogy. It is the same theorem
applied to a different conserved quantity.

\end{proof}

\subsection*{2.3 The Wave Equation for \ensuremath{\psi}}

\begin{theorem}[Polarization Wave Equation]
The polarization field \ensuremath{\psi} satisfies
the discrete wave equation on the FCC lattice:
\end{theorem}

\begin{equation*}
\partial ^{2}\psi_{i} / \partial \tau^{2} = c^{2}_{\mathrm{TMS}} \Delta_{\mathrm{FCC}} \psi_{i}
\end{equation*}

with propagation speed \ensuremath{c_{\mathrm{TMS}}} = \ensuremath{1/\sqrt{2}} lattice spacings per flop. In the
continuum limit:

\begin{equation*}
\partial ^{2}\psi / \partial t^{2} - c^{2}\nabla ^{2}\psi = 0
\end{equation*}

\begin{proof}
The propagation-speed argument of Paper 1 Section 6.5 depends
only on the FCC lattice geometry --- specifically, on the second-order
tensor \ensuremath{\Sigma_{j}} \ensuremath{r_{j}\mu} \ensuremath{r_{j}\nu} = \ensuremath{4\ell^{2}_{p}} \ensuremath{\delta \mu \nu} computed over the 12 nearest neighbors.
This tensor is a property of the lattice, not of the field being
propagated. The same argument applies to any scalar field that
propagates by nearest-neighbor coupling on the FCC lattice. Therefore
\ensuremath{c_{\mathrm{TMS}}} = \ensuremath{1/\sqrt{2}} for \ensuremath{\psi .} The continuum limit follows by the Laplacian
convergence argument of Paper 1 Section 6.6, again depending only on
lattice geometry.

\end{proof}

\begin{corollary}[No New Parameter]
The polarization field introduces no new
propagation speed, no new coupling constant, and no new length scale. It
inherits \ensuremath{\ell_{p},} \ensuremath{t_{p},} and c from Paper 1 without modification.
\end{corollary}

\subsection*{2.4 Independence of \ensuremath{\varphi} and \ensuremath{\psi} in the Bulk}

The confirmation field and the polarization field share a substrate but
propagate independently in the bulk. A wave in \ensuremath{\varphi} describes how
confirmation density propagates regardless of binary content; a wave in
\ensuremath{\psi} describes how polarization content propagates regardless of
confirmation density. The two fields couple only at confirmation events
themselves --- the moments at which a new bit is added to the causal
record. The coupling mechanism is the subject of Section III.

\section*{III. Confirmation Events: The Coupling of \ensuremath{\varphi} and \ensuremath{\psi}}
\addcontentsline{toc}{section}{III. Confirmation Events: The Coupling of \ensuremath{\varphi} and \ensuremath{\psi}}

\subsection*{3.1 The Atomicity of the Flop Revisited}

Paper 1 Section 1.5 established that a flop is indivisible: the locking
of a bit, the firing of 12 outgoing signals, the consumption of those
signals as 4 synthetic bit positions, and the Z-displacement to the next
layer all occur as a single act with no sub-flop time. Paper 2 examines
the substructure of this single act with respect to the polarization
landscape --- not by subdividing the flop temporally (which is
forbidden) but by identifying the geometric components present at the
instant the flop occurs.

At the moment a confirmation event occurs at synthetic bit position x in
layer \ensuremath{S_{n+1}:}

\begin{itemize}
\item
  Three touch vector signals are arriving at x, one from each of the
  three confirming agents in the parent layer \ensuremath{S_{n}.} Each signal carries
  the polarization of its source confirmation: \ensuremath{\chi_{1},} \ensuremath{\chi_{2},} \ensuremath{\chi_{3}} \ensuremath{\in} \{\ensuremath{-1,} +1\}.
\item
  \ensuremath{U_{\mathrm{sh}}} contributes its standard unbiased stochastic input --- the same
  maximum-entropy field that supplies latent bits across the substrate
  (Axiom 0; Paper 1 Section 1.4).
\item
  The triadic agreement rule fires, locking the synthetic bit position
  into a confirmed state and selecting its binary value \ensuremath{\bar{B}} \ensuremath{\in} \{0, 1\}.
\end{itemize}

These components are not temporally sequential --- they are the
geometric content of the single indivisible act.

\subsection*{3.2 The Local Polarization Amplitude}

\begin{definition}[Local Polarization Amplitude]
At a synthetic bit position x
in layer \ensuremath{S_{n+1},} the local polarization amplitude is the average of the
polarizations carried by the three incoming touch vector signals from
x's parent agents in \ensuremath{S_{n}:}
\end{definition}

\begin{equation*}
\Pi_{\psi}(x) = (\chi_{1} + \chi_{2} + \chi_{3}) / 3
\end{equation*}

\ensuremath{\Pi_{\psi}(x)} ranges over the four discrete values \{\ensuremath{-1,} \ensuremath{-1/3,} +1/3, +1\},
corresponding to the four possible parent configurations under the
symmetry of the three input slots.

\ensuremath{\Pi_{\psi}(x)} is the same \ensuremath{\Pi_{\psi}} that propagates as the conjugate excitation of
\ensuremath{\psi} in Section II --- measured at the specific synthetic-bit position x.
The bulk wave equation describes how \ensuremath{\Pi_{\psi}} evolves across the lattice
between confirmation events; the local polarization amplitude at x is
its value at the moment a confirmation occurs there.

\subsection*{3.3 Synthetic Bits Carry Position, Not Value}

\begin{proposition}
A synthetic bit, prior to its confirmation event, is a
geometric appointment --- a position in \ensuremath{S_{n+1}} at which a confirmation
must occur --- and does not carry a binary value of its own. Its value
is determined at the moment of confirmation by the joint action of
incoming polarization and \ensuremath{U_{\mathrm{sh}}.}
\end{proposition}

\begin{proof}
By the \ensuremath{1\to 4} multiplier (Paper 1 Section 4.2), each primary
confirmation produces 4 synthetic bit positions through the geometric
consumption of its 12 outgoing signals. The geometric act that produces
a synthetic bit position is the projection of touch vectors forward
through the FCC lattice; this projection determines location but does
not by itself determine binary value. Conferring a value on the
synthetic bit at the moment of its production would require selecting
between Possibility 1 (inherited value rule, violating Generative
Economy) and Possibility 2 (independent randomness at production,
violating Axiom 3) from Section 1.2. Possibility 3 --- biased stochastic
resolution at the moment of confirmation --- is the unique geometrically
consistent option.

\end{proof}

This distinction between geometric appointment and valued outcome
resolves an apparent ambiguity in the substrate's dynamics: synthetic
bits and confirmed bits are the same physical event seen at
production-time and stabilization-time respectively. The \ensuremath{1\to 4} multiplier
produces appointments; the confirmation event produces values.

\section*{IV. The Born-Rule Structural Correspondence}
\addcontentsline{toc}{section}{IV. The Born-Rule Structural Correspondence}

\subsection*{4.1 The Resolution Probability}

\begin{theorem}[Polarization-Biased Resolution]
Given a synthetic bit position
x with local polarization amplitude \ensuremath{\Pi_{\psi}(x)} \ensuremath{\in} \ensuremath{[-1,} +1], the
probability that the confirmation event at x resolves as \ensuremath{\bar{B}} = 1 is:
\end{theorem}

\emph{\ensuremath{P(\bar{B}} = 1 \textbar{} \ensuremath{\Pi_{\psi})} = (1 + \ensuremath{\Pi_{\psi})} / 2}

and equivalently the probability of \ensuremath{\bar{B}} = 0 is:

\emph{\ensuremath{P(\bar{B}} = 0 \textbar{} \ensuremath{\Pi_{\psi})} = (1 \ensuremath{-} \ensuremath{\Pi_{\psi})} / 2}

These probabilities sum to unity by construction and are non-negative
for all \ensuremath{\Pi_{\psi}} \ensuremath{\in} \ensuremath{[-1,} +1].

\begin{proof}
The resolution probability must satisfy four constraints, each
forced by a primitive of the TMS:

\begin{itemize}
\item
  \ensuremath{\Pi_{\psi}} = +1 (all parents confirmed \ensuremath{\bar{B}} = 1) \ensuremath{\to} \ensuremath{P(\bar{B}} = 1) = 1. Required by
  Axiom 3: if all three incoming signals carry \ensuremath{\chi} = +1, dissipating any
  of that polarization through a non-certain resolution would violate
  signal conservation in the polarization content.
\item
  \ensuremath{\Pi_{\psi}} = \ensuremath{-1} (all parents confirmed \ensuremath{\bar{B}} = 0) \ensuremath{\to} \ensuremath{P(\bar{B}} = 1) = 0. Required by
  the same argument with the opposite polarization.
\item
  \ensuremath{\Pi_{\psi}} = 0 (no net bias) \ensuremath{\to} \ensuremath{P(\bar{B}} = 1) = 1/2. Required by the
  maximum-entropy character of \ensuremath{U_{\mathrm{sh}}} (Axiom 0): in the absence of bias,
  the stochastic kick from \ensuremath{U_{\mathrm{sh}}} resolves the synthetic bit with equal
  probability between its two possible values.
\item
  Monotonicity and continuity in \ensuremath{\Pi_{\psi}.} Required by the Principle of
  Generative Economy: any non-monotonic or discontinuous dependence on
  \ensuremath{\Pi_{\psi}} would introduce structure beyond what is forced by (i)--(iii),
  violating the minimum-sufficient selection criterion.
\end{itemize}

The unique function satisfying (i), (ii), (iii), and (iv) is the linear
map \ensuremath{P(\bar{B}} = 1 \textbar{} \ensuremath{\Pi_{\psi})} = (1 + \ensuremath{\Pi_{\psi})/2.} No alternative mapping
satisfies all four constraints without an additional parameter (such as
a non-linear curve shape).

\end{proof}

\subsection*{4.2 Connection to the Born Rule}

The resolution probability \ensuremath{P(\bar{B}} = 1 \textbar{} \ensuremath{\Pi_{\psi})} = (1 + \ensuremath{\Pi_{\psi})/2} is
structurally analogous to the Born rule for projective measurement in
quantum mechanics (Born, 1926). The correspondence consists of three
parallel components:

\begin{itemize}
\item
  Unitary propagation between measurements \ensuremath{\leftrightarrow} wave propagation between
  confirmation events. In QM, the state evolves deterministically via
  the Schrödinger equation between measurements. In TMS, the
  polarization field \ensuremath{\psi} and its excitation \ensuremath{\Pi_{\psi}} evolve deterministically
  via the FCC wave equation between confirmation events.
\item
  Measurement outcome biased by amplitude \ensuremath{\leftrightarrow} resolution outcome biased by
  \ensuremath{\Pi_{\psi}.} In QM, the outcome of a projective measurement is stochastic but
  its probability is determined by the local wave amplitude. In TMS, the
  outcome of a confirmation event is stochastic but its probability is
  determined by the local polarization amplitude.
\item
  Born probability rule \ensuremath{\leftrightarrow} linear polarization-to-probability mapping. In
  QM, P(outcome) = \textbar\ensuremath{\langle \text{outcome}} \textbar{} \ensuremath{\psi \rangle}\textbar\ensuremath{^{2}} --- a
  quadratic mapping from complex amplitude to probability. In TMS, \ensuremath{P(\bar{B}} =
  1) = (1 + \ensuremath{\Pi_{\psi})/2} --- a linear mapping from real polarization to
  probability.
\end{itemize}

Status of the Correspondence. The mapping (1 + \ensuremath{\Pi_{\psi})/2} is the Born rule
restricted to two-outcome projective measurement on real-valued
amplitudes with no phase information. The recovery of the full quadratic
\textbar\ensuremath{\cdot}\textbar\ensuremath{^{2}} mapping requires complex amplitudes and phase
coherence --- structures that the TMS does not yet exhibit at the level
of the polarization field \ensuremath{\psi} alone. Phase information in the TMS is
carried by the time-evolution of the \ensuremath{(\psi ,} \ensuremath{\Pi_{\psi})} pair across many flops,
not by an instantaneous complex amplitude at any single lattice site.
The derivation of the full complex Born rule from the macroscopic
interference behavior of the polarization wave equation is reserved for
subsequent work. The present claim is one of structural correspondence
--- the same skeleton of ``deterministic wave evolution +
amplitude-biased projective measurement'' --- exactly parallel in status
to the [[4,2,2]] structural correspondence of Paper 1, where the
stabilizer algebra is faithfully reproduced but the full quantum
amplitude formalism is deferred. \hfill\ensuremath{\square}

\begin{remark}
This identification dissolves the standard interpretive
distinction between unitary evolution and measurement collapse. In the
TMS, both are aspects of the same substrate dynamics. The wave equation
governs propagation; the confirmation event resolves the value; \ensuremath{U_{\mathrm{sh}}}
supplies the irreducible stochasticity. No additional postulate is
required to separate the two regimes. The ``collapse'' of standard
quantum mechanics is, in TMS terms, the per-flop projection of
polarization amplitude onto the binary alphabet.
\end{remark}

\subsection*{4.3 Positioning Against Parallel Programs}

Several active programs aim to derive quantum measurement from
substrate-level dynamics, and the polarization-biased resolution rule of
\S\,4.1 should be situated explicitly against the two most directly
comparable proposals.

't Hooft's Cellular Automaton Interpretation of quantum mechanics ('t
Hooft, 2016) proposes that the apparent randomness of quantum
measurement is the consequence of observing a deterministic underlying
cellular automaton in the ``wrong basis'' --- a basis in which the
deterministic dynamics appear as superposition and the deterministic
resolution appears as stochastic collapse. The CAI's deterministic
ontology is realized in TMS in a partial sense: the polarization-biased
resolution rule fixes the probability of each outcome from the local
polarization amplitude, so the substrate is deterministic in its bias
but stochastic in its residual selection. The residual stochasticity
comes from \ensuremath{U_{\mathrm{sh}},} which by Axiom 0 is algorithmically incompressible ---
not a hidden deterministic dynamic but an irreducible source. TMS is
therefore neither a hidden-variable theory in the strict CAI sense nor a
Copenhagen-style ontological collapse theory; it is a third option in
which the bias is geometric and the residual randomness is forced by the
maximum-entropy substrate.

Gorard's derivation of quantum measurement in the Wolfram model (Gorard,
2020) proceeds through the multiway evolution graph: a single substrate
state evolves into many possible successors under non-confluent
hypergraph rewriting, and observers impose effective causal invariance
via a Knuth-Bendix completion that collapses distinct multiway branches
into a single observed thread. TMS achieves a similar outcome through a
different mechanism. There is no multiway evolution graph --- every flop
produces a single confirmed configuration on a single FCC lattice,
determined by the polarization-biased resolution rule with \ensuremath{U_{\mathrm{sh}}} supplying
the stochastic component. The substrate of TMS is not a superposition of
histories projected onto a single observed branch; it is a single
history whose per-event outcomes are stochastic by the maximum-entropy
properties of \ensuremath{U_{\mathrm{sh}}.} The two pictures converge on the operational
appearance of quantum measurement --- wave-like evolution between
events, probabilistic resolution at each event --- while disagreeing on
whether the substrate carries one history or many.

The decisive structural difference common to both contrasts is that TMS
derives the linear amplitude-to-probability mapping (1 + \ensuremath{\Pi_{\psi})/2} from
signal conservation and the maximum-entropy character of \ensuremath{U_{\mathrm{sh}}} alone,
without invoking either a hidden deterministic basis ('t Hooft) or a
multiway branching structure (Gorard). The mapping is the unique
function satisfying the four geometric constraints of \S\,4.1, and its
emergence requires no additional ontological machinery beyond what was
already established in Paper 1.

\section*{V. Universal Computation from Triadic Confirmation}
\addcontentsline{toc}{section}{V. Universal Computation from Triadic Confirmation}

\subsection*{5.1 The 3-Input Majority Gate}

\begin{theorem}[Majority Gate Emergence]
Every confirmation event in the TMS
implements a probabilistic 3-input majority gate over the polarizations
of its parent confirmations, with the bias toward the deterministic
majority outcome set by the magnitude of the polarization amplitude
\textbar \ensuremath{\Pi_{\psi}}\textbar.
\end{theorem}

\begin{proof}
From Section 3.2, the local polarization amplitude at synthetic
bit position x is the average of three parent polarizations: \ensuremath{\Pi_{\psi}(x)} =
\ensuremath{(\chi_{1}} + \ensuremath{\chi_{2}} + \ensuremath{\chi_{3})/3.} The four possible parent configurations (up to
permutation) and their resolution probabilities are:

\begin{itemize}
\item
  (+1, +1, +1): \ensuremath{\Pi_{\psi}} = +1, \ensuremath{P(\bar{B}} = 1) = 1, \ensuremath{P(\bar{B}} = 0) = 0.
\item
  (+1, +1, \ensuremath{-1):} \ensuremath{\Pi_{\psi}} = +1/3, \ensuremath{P(\bar{B}} = 1) = 2/3, \ensuremath{P(\bar{B}} = 0) = 1/3.
\item
  (+1, \ensuremath{-1,} \ensuremath{-1):} \ensuremath{\Pi_{\psi}} = \ensuremath{-1/3,} \ensuremath{P(\bar{B}} = 1) = 1/3, \ensuremath{P(\bar{B}} = 0) = 2/3.
\item
  \ensuremath{(-1,} \ensuremath{-1,} \ensuremath{-1):} \ensuremath{\Pi_{\psi}} = \ensuremath{-1,} \ensuremath{P(\bar{B}} = 1) = 0, \ensuremath{P(\bar{B}} = 0) = 1.
\end{itemize}

In the deterministic limit \textbar \ensuremath{\Pi_{\psi}}\textbar{} \ensuremath{\to} 1 (achieved when
all three parents agree), the resolution is certain to match the
parents. In the intermediate cases, the resolution is biased toward the
majority value with probability 2/3. This is the canonical probabilistic
majority function on three binary inputs.

\end{proof}

\subsection*{5.2 Stabilizer Repair and Effective Determinism}

The 1/3 probability of ``minority'' resolutions in the intermediate
cases is not noise to be ignored. It is a single-bit error rate. By
Paper 1 Section 5, every tetrahedral confirmation unit detects and
repairs single-bit errors within the same flop via the [[4,2,2]]
stabilizer dynamics. The synthetic bits produced at the current flop
repair the missing or anomalous confirmation, restoring stabilizer
parity.

\begin{theorem}[Effective Determinism at Scale]
When the [[4,2,2]]
stabilizer repair mechanism is applied across a connected region of the
lattice in which polarizations are coherent, the effective behavior of
the substrate at scale is deterministic 3-input majority computation.
\end{theorem}

Proof sketch. The single-bit error rate per confirmation in the
intermediate cases is 1/3. The [[4,2,2]] code has distance 2 ---
any single-bit error is detected, and the synthetic repair process
applied at the next flop projects the local configuration toward the
nearest stabilizer-satisfying state. In a coherent region (where the
polarization landscape varies slowly on the scale of single confirmation
events), the nearest stabilizer-satisfying state is the deterministic
majority result. Therefore the effective per-confirmation error rate
after stabilizer projection decays geometrically with depth into the
coherent region. At causal depth d, the effective probability of an
erroneous resolution is bounded above by \ensuremath{(1/3)^{d},} which vanishes
rapidly. The substrate's macroscopic behavior is therefore effectively
deterministic majority computation in any coherent region. A complete
derivation of the decay rate and the exact dependence on
coherence-region geometry is provided in Paper 3 \S\,2.5. \hfill\ensuremath{\square}

\subsection*{5.3 Universality}

\begin{corollary}[Computational Universality]
The TMS substrate supports
universal classical computation.
\end{corollary}

\begin{proof}
The 3-input majority gate is universal in combination with a
constant input. Specifically, MAJ(x, y, 0) = AND(x, y), MAJ(x, y, 1) =
OR(x, y), and NOT can be implemented by polarization inversion --- a
physical reflection of the polarization landscape across the \ensuremath{\chi} = 0 axis,
achievable in TMS by routing through a confirmation region whose parent
polarizations are uniformly opposite to the input. The set \{AND, OR,
NOT\} is functionally complete (Shannon, 1948).

\end{proof}

The substrate is therefore a universal computer by geometric necessity.
No additional design is required --- the triadic confirmation rule,
applied to a binary alphabet on a fault-tolerant FCC lattice, is
sufficient.

\section*{VI. Programs, Selection, and Computational Irreducibility}
\addcontentsline{toc}{section}{VI. Programs, Selection, and Computational Irreducibility}

\subsection*{6.1 The Flop Operator and Configurations}

\begin{definition}[Flop Operator]
The flop operator T maps the polarization
landscape at layer \ensuremath{S_{n}} to the (probabilistic) polarization landscape at
layer \ensuremath{S_{n+1}:}
\end{definition}

\begin{equation*}
T: \psi (S_{n}) \mapsto \psi (S_{n+1})
\end{equation*}

T is defined by the local resolution rule of Section IV combined with
the global geometric projection of the \ensuremath{1\to 4} multiplier from \ensuremath{S_{n}} into \ensuremath{S_{n+1}.}
T is probabilistic in general (due to the \ensuremath{U_{\mathrm{sh}}} contribution) and becomes
effectively deterministic in coherent regions (per Section 5.2).

\begin{definition}[Configuration]
A configuration C is a finite spatial
pattern of confirmed polarizations across some bounded region of the
lattice extending across a causal depth window d:
\end{definition}

\emph{C = \{\ensuremath{\chi (x,} \ensuremath{\tau )} : x \ensuremath{\in} R, \ensuremath{\tau} \ensuremath{\in} \ensuremath{[\tau_{0},} \ensuremath{\tau_{0}} + d]\}}

where R is the spatial support of the configuration.

\subsection*{6.2 Programs}

\begin{definition}[Program]
A configuration C is a program of period n iff:
\end{definition}

\begin{itemize}
\item
  Causal closure under T: \ensuremath{T^{n}[C]} = C with probability approaching
  1 in the limit of large coherent surround. The configuration's
  polarization landscape reproduces itself, possibly up to lattice
  translation, after n flops.
\item
  Algorithmic compressibility: K(C) \textless{} \textbar C\textbar.
  There exists a description of C shorter than its raw enumeration.
\end{itemize}

Programs of period 1 are static patterns. Programs of period n
\textgreater{} 1 are oscillators or propagating patterns. A program of
period n that translates by some lattice vector v under \ensuremath{T^{n}} is a
propagating program --- the TMS analogue of a particle, though the
connection to physical matter is deferred to subsequent work.

\begin{remark}
Condition (b) is the formal criterion that distinguishes a
program from a coincidence in \ensuremath{U_{\mathrm{sh}}.} Random configurations drawn directly
from the maximum-entropy substrate satisfy K(C) \ensuremath{\approx} \textbar C\textbar{}
almost surely (Kolmogorov, 1965; Axiom 0). Only configurations with K(C)
\textless{} \textbar C\textbar{} have structure beyond noise --- beyond
what \ensuremath{U_{\mathrm{sh}}} would produce by accident. Programs in the TMS are not
configurations that ``look organized'' by inspection; they are
configurations that are organized in the strict Kolmogorov sense.
\end{remark}

\begin{remark}[Status of K(C) in the Program Definition]
The
Kolmogorov complexity K(C) is not computable in general (Kolmogorov,
1965). The definition of a program via K(C) \textless{}
\textbar C\textbar{} therefore raises an apparent question: how does the
substrate select program-bearing configurations if it cannot compute the
criterion that defines them?
\end{remark}

The answer is that the substrate does not compute K(C). K(C) appears
here as a property of configurations selected by the substrate's
geometric dynamics, not as a criterion the substrate evaluates and
applies. The selection mechanism is the stabilizer dynamics plus
T-closure plus the coherence threshold --- all of which are local
geometric operations on the polarization landscape, computable at every
flop by the substrate's own update rule. Configurations that survive
this selection happen to satisfy K(C) \textless{} \textbar C\textbar{}
as a consequence of two structural facts: (a) random samples drawn
directly from \ensuremath{U_{\mathrm{sh}}} have K(C) \ensuremath{\approx} \textbar C\textbar{} almost surely by
Kolmogorov's incompressibility theorem; (b) configurations that survive
geometric selection are by definition not random samples from \ensuremath{U_{\mathrm{sh}},} since
they exhibit the structural regularities required by T-closure and
stabilizer parity.

The implication runs in one direction only. Geometric selection implies
K(C) \textless{} \textbar C\textbar{} as an emergent property of the
survivors. K(C) \textless{} \textbar C\textbar{} does not need to be
checked, computed, or enforced --- it is what the selected
configurations turn out to have. The substrate operates entirely on
local geometric criteria; the algorithmic-compressibility
characterization is the external observer's description of what those
criteria pick out.

This distinction matters for the exhaustion analysis of Paper 4 \S\,VI.
Exhaustion is not a substrate event triggered by the substrate
``noticing'' that novelty has run out; it is a property of the
configuration space that emerges from the finite size of \ensuremath{\Sigma (R).} The
substrate continues operating throughout. The transition to
reinterpretation at \ensuremath{\tau} \textgreater{} \ensuremath{\tau *(R)} is forced by the structural
situation (Axiom 4 forbids overwriting the confirmed bits; Axiom 3
forbids losing the signals that produced them; direct sampling extracts
no further novelty from a saturated configuration space) --- not by the
substrate computing anything about K(C).

\subsection*{6.3 The Selection Mechanism}

The TMS provides an automatic selection mechanism that separates
programs from noise without any external evaluator.

\begin{theorem}[Stabilizer Selection]
Configurations that satisfy
[[4,2,2]] stabilizer parity and are causally closed under T
persist across flops with probability approaching 1 in the limit of
large coherent surround. Configurations that violate stabilizer parity
are repaired toward the nearest parity-satisfying configuration.
Configurations that are not causally closed dissolve into \ensuremath{U_{\mathrm{sh}}-\text{like}} noise
across O(d) flops, where d is the causal depth across which they would
need to maintain coherence.
\end{theorem}

Proof sketch. By Paper 1 Section 5, single-bit deviations from any
locally stabilizer-satisfying configuration are detected and repaired
within the same flop. Configurations that satisfy stabilizer parity are
local fixed points of the repair process. Configurations that violate
parity are projected by repair toward the nearest stabilizer-satisfying
state. For a configuration to persist as a program, it must additionally
close under T --- i.e., the natural forward evolution of its
polarization landscape must reproduce its current state. Configurations
that satisfy stabilizer parity but not causal closure under T are
repaired in place but drift under wave dynamics, dissolving across O(d)
flops where d is the spatial extent of the configuration relative to the
wave propagation speed \ensuremath{c_{\mathrm{TMS}}} = \ensuremath{1/\sqrt{2}.} A complete derivation of the
dissolution rate is provided in Paper 3 \S\,2.5. \hfill\ensuremath{\square}

\begin{remark}
This is selection without natural selection. The TMS substrate
distinguishes programs from noise by the same mechanism it distinguishes
confirmed states from unconfirmed states: stabilizer projection. There
is no separate evaluator, no fitness function, no environmental judge.
The geometry is the selector. Configurations that compute persist;
configurations that do not compute dissolve.
\end{remark}

\subsection*{6.4 Computational Irreducibility of Programs}

The arc from Axiom 0 to programs is a continuation of the arc that Paper
1 traced from Axiom 0 to the GUP. A program of causal depth d is a
configuration whose existence requires d flops of confirmation to be
specified. By Axiom 0, the \ensuremath{U_{\mathrm{sh}}} substrate from which each layer's
contribution to the program is drawn is algorithmically incompressible:
K(n) \ensuremath{\approx} n.~The down-vector chain encoding a program of depth d therefore
has a lower bound on its compressibility set by the \ensuremath{U_{\mathrm{sh}}} contributions at
each layer.

For programs (configurations with K(C) \textless{} \textbar C\textbar),
this lower bound is non-trivial but non-vanishing: programs compress
because of their internal structure, but they do not compress
arbitrarily because their causal history records irreducible
contributions from \ensuremath{U_{\mathrm{sh}}} at every layer.

\begin{proposition}[Program Irreducibility]
A program of causal depth d has
minimum description length bounded below by \ensuremath{K_{\psi}(d)} --- the algorithmic
complexity of its polarization history along the depth axis. This bound
is strictly positive for all programs and grows monotonically with d for
non-trivial programs.
\end{proposition}

This bound is the computational content of the \ensuremath{\beta} = f(d, \ensuremath{K_{\psi})} prediction
sharpened in Paper 1 Section 7.2. The function f, refined to its full
functional dependences in Paper 3 \S\,6.2, has its physical meaning fixed
by Paper 2: \ensuremath{K_{\psi}} is the algorithmic complexity of the program running
through the bit's causal history, and the GUP correction at the bit's
position is set by the irreducibility of that program. Position
measurement of a deeply embedded bit requires reading out its program,
which by definition cannot be compressed below \ensuremath{K_{\psi}(d).}

\subsection*{6.5 The Role of Volume 1, Paper 3}

The selection mechanism described in Section 6.3 is the substrate-level
prerequisite for learning, hierarchy, and the emergence of subsystems
with internal organization. Paper 3 develops the macroscopic
consequences: how programs of increasing causal depth come to dominate
the polarization landscape over the long flop-time history of a stable
cascade instance, how the substrate concentrates computation into
coherent regions through the bistability mechanism established there,
how nested programs (programs running inside the spatial support of
other programs) produce hierarchical structure, and how the first
subsystem --- the smallest configuration that transforms patterns rather
than merely persisting them --- emerges from the same axioms. The
dissolution rate of non-closed configurations and the effective
error-rate decay with coherence depth are both closed in Paper 3 \S\,2.5 by
derivations consistent with the bounds noted above.

\section*{VII. The PGE-Forced Wave-Form of Higher-Order Reinterpretation}
\addcontentsline{toc}{section}{VII. The PGE-Forced Wave-Form of Higher-Order Reinterpretation}

The substrate's first-order dynamics are wave propagation --- the wave
equation derived in Paper 1 \S\,VI on the FCC confirmation field is the
substrate-level mechanism by which information propagates across the
agent manifold. The structurally prior question we address here is: when
reinterpretation at any hierarchical level introduces a higher-order
dynamical phase, what form does that higher-order dynamics take, and
why?

We show that higher-order reinterpretation is \emph{wave-form by
structural necessity}, not by additional axiom. The argument is four
steps under PGE. The theorem developed here is invoked by Paper 5 \S\,III
to resolve the apparent tension between the first-order substrate wave
equation and the second-order wave structure that reinterpretation
produces.

\subsection*{7.1 The Four-Step Argument}

\begin{theorem}[PGE-Forced Wave-Form of Reinterpretation]
The
Principle-of-Generative-Economy-selected continuation of the substrate
at any reinterpretation boundary takes wave-form dynamics.
\end{theorem}

\begin{proof}
Four steps.

\emph{(i) First-order substrate behavior is wave propagation.} By Paper
1 \S\,VI, the confirmation field \ensuremath{\varphi} on the FCC agent manifold satisfies the
wave equation \ensuremath{\partial ^{2}\varphi /\partial t^{2}} \ensuremath{-} \ensuremath{c^{2}\nabla ^{2}\varphi} = 0 in the continuum limit, with c = \ensuremath{c_{\mathrm{TMS}}}
= \ensuremath{1/\sqrt{2}} lattice spacings per flop. Wave propagation is the substrate's
native vocabulary for moving information across the agent manifold.

\emph{(ii) PGE forbids importing mechanisms outside the substrate's
vocabulary.} The Principle of Generative Economy, as developed in Paper
1 \S\,1.3 and applied throughout this volume, selects the minimum
sufficient closure at every scale. A continuation that requires a new
propagation mechanism --- one that does not already operate at the
substrate level --- adds structure beyond what the existing geometry
forces. PGE excludes this. Any higher-order dynamics that PGE selects
must reuse the substrate's existing propagation vocabulary.

\emph{(iii) The substrate's only propagation vocabulary is wave
propagation.} The FCC geometry, the triadic confirmation rule, and the
[[4,2,2]] stabilizer structure collectively determine the
substrate's repertoire of information-moving operations. The wave
equation is the unique continuum-limit propagation rule that this
geometry supports --- first established at substrate level in Paper 1,
and shown to be form-invariant under recursive coarse-graining (G
operator) in \S\,IV of this paper. There is no second native propagation
mechanism. The substrate's vocabulary contains exactly one verb for
moving information through agent space: the wave.

\emph{(iv) Therefore higher-order reinterpretation is wave-form.}
Combining (ii) and (iii): since PGE forbids importing new mechanisms and
the only available mechanism is wave propagation, the PGE-selected
continuation at a reinterpretation boundary uses wave propagation at the
new scale. The recursion is therefore a \emph{wave hierarchy}. Each
order of wave reads the structure left by the previous orders,
propagating in the level-k flop time at the level-k propagation speed
\ensuremath{c^{(k)}.}

\end{proof}

\textbf{Status of this result.} \emph{Derived.} The four-step argument
follows from the wave equation of Paper 1 \S\,VI, the PGE statement of
Paper 1 \S\,1.3, and the form-invariance of the wave equation under G that
we established earlier in this paper. No additional structure is
introduced.

\subsection*{7.2 Quantization and the Erosion Picture (Previews of Paper 5 \S\,3.5)}

\textbf{Remark on quantization (preview of Paper 5 \S\,3.5).} A wave
hierarchy with closure constraints at each level forces a specific
mechanism for stable excitations: packets must settle into
configurations where the wave geometry permits standing-wave closure ---
wave troughs in the accumulated wake of prior propagation. Quantization
is not a separate axiom; it is what falls out of having a recursive wave
hierarchy at all. This is developed structurally in Paper 5 \S\,3.5 and
applied to field content in Paper 5 \S\,IV.

\textbf{Remark on the erosion picture.} The recursive wave hierarchy is,
equivalently, an accumulating wake: every prior propagation event leaves
a residue in the substrate's effective topography, and new propagation
events ride atop the topography that prior events have already carved.
At any given moment, the substrate's working geometry is a superposition
of an ancient layer (deep channels carved by very early dynamics,
effectively rigid at any subsystem-accessible timescale), an
intermediate layer (ongoing dynamics across the current epoch), and a
transient layer (flash-flood channels carved by recent events that have
not yet relaxed). This picture is developed for cosmological purposes in
Paper 5 \S\,3.5 and \S\,V; the bare structural claim --- that field-effective
stability is an erosion property of the wave hierarchy rather than a
fundamental axiom --- is what the four-step theorem above establishes.

\section*{VIII. Predictions and Directions for Subsequent Work}
\addcontentsline{toc}{section}{VIII. Predictions and Directions for Subsequent Work}

\subsection*{8.1 Empirical Signatures of Polarization Refraction}

Although the polarization field \ensuremath{\psi} propagates at \ensuremath{c_{\mathrm{TMS}}} = \ensuremath{1/\sqrt{2}} across the
bulk lattice (Section 2.3) --- the same speed as \ensuremath{\varphi} --- its interaction
with the confirmation density at the leading edge of a cascade can
produce apparent variations in effective propagation velocity. In
regions where \ensuremath{\varphi} is rapidly varying (sharp confirmation gradients) and
the polarization landscape is structured, coherent interference between
the \ensuremath{\varphi} and \ensuremath{\psi} wave fields produces what an external observer would
interpret as a position-dependent index of refraction.

This is not a modification of the underlying wave speed. It is
interference between two independent fields on the same lattice,
observed at coarse-grained scale. The phenomenon is parallel to real
optical refraction in standard physics, where the underlying
electromagnetic field travels at c everywhere and the apparent slowdown
in a dielectric medium is interference with the medium's polarization
response. The TMS provides a geometric microstructure for this same
phenomenon.

A specific prediction: the apparent variation in effective propagation
velocity scales with the local polarization gradient
\textbar\ensuremath{\nabla \psi}\textbar{} and the local confirmation density \ensuremath{\varphi .} Regions with
steep \ensuremath{\psi -\text{gradients}} (sharp boundaries between coherent polarization
domains) act as effective interfaces. Regions with shallow \ensuremath{\psi -\text{gradients}}
propagate signals at the bare lattice speed. The functional form is
derivable from the coupled wave equations of \ensuremath{(\varphi ,} \ensuremath{\psi )} and is developed
further in Paper 3 alongside the related dissolution-rate and
error-decay results.

\subsection*{8.2 The Sharpened GUP Prediction}

Paper 1 Section 7.2 introduced the prediction \ensuremath{\beta} = f(d, \ensuremath{K_{\psi}).} Paper 2
has filled in the meaning of \ensuremath{K_{\psi}:} it is the Kolmogorov complexity of
the program running through the bit's causal history, where ``program''
has the formal definition of Section 6.2. The prediction therefore
acquires its computational content here. Bits embedded in deep,
computationally rich programs (high \ensuremath{K_{\psi}} at high d) exhibit measurably
larger GUP corrections than bits in shallow, computationally trivial
configurations at comparable energy scales.

This is a falsifiable prediction with two empirical correlates. First,
quantum systems with rich internal computational structure (highly
entangled, deeply prepared states) should exhibit measurably larger
minimum-uncertainty floors than systems with trivial internal structure
at the same total energy. Second, the GUP correction should not be a
fixed property of a particle species but should vary with the particle's
preparation history. Both are accessible in principle to Planck-scale
phenomenology programs (Hossenfelder, 2013) and to next-generation
quantum optomechanics experiments. Paper 3 \S\,6.2 further refines f to
account for surround-imprinting and subsystem context, broadening the
empirical correlates further.

\subsection*{8.3 The Born Rule from Coarse-Graining}

The full quadratic Born rule of standard quantum mechanics is a target
for subsequent work. The structural correspondence of Section IV
establishes a linear amplitude-to-probability mapping for two-outcome
projective measurement with real polarization amplitudes. The recovery
of the quadratic \textbar\ensuremath{\cdot}\textbar\ensuremath{^{2}} mapping requires complex amplitudes,
which in the TMS are expected to arise from the phase structure of
coupled \ensuremath{(\psi ,} \ensuremath{\Pi_{\psi})} propagation across many flops. The precise derivation
is deferred. The key conceptual claim has already been established: the
substrate provides a single mechanism --- polarization-biased
confirmation against \ensuremath{U_{\mathrm{sh}}} --- that unifies what standard quantum
mechanics treats as two separate postulates (unitary evolution and
measurement collapse).

\subsection*{8.4 Programs, Particles, and Subsequent Volumes}

The propagating programs of Section 6.2 --- periodic polarization
configurations that translate under \ensuremath{T^{n}} --- are the substrate-level
structures from which physical matter is expected to emerge in
subsequent volumes. The connection is anticipated rather than asserted
here. The substrate-level claim is independently meaningful: the TMS
supports a discrete spectrum of stable propagating computational
patterns. Whether and how these correspond to the particle spectrum of
the Standard Model is the subject of later work in Volume 1 and Volume
2.

\section*{Conclusion: A Complete Computational Substrate}
\addcontentsline{toc}{section}{Conclusion: A Complete Computational Substrate}

Building on the geometric and ontological architecture of Paper 1, Paper
2 has shown that the TMS substrate supports computation by geometric
necessity, with no additional axioms or free parameters. The following
results emerge from the primitives of Paper 1 alone:

\begin{itemize}
\item
  The polarization field \ensuremath{\psi} propagates by a homogeneous wave equation on
  the FCC lattice at \ensuremath{c_{\mathrm{TMS}}} = \ensuremath{1/\sqrt{2},} identical in form to the confirmation
  field \ensuremath{\varphi} --- the propagation speed and lattice structure are inherited,
  not re-derived.
\item
  The conjugate pair \ensuremath{(\psi ,} \ensuremath{\Pi_{\psi})} emerges from the same conservation
  argument that produced \ensuremath{(\varphi ,} \ensuremath{\Pi )} in Paper 1, applied to polarization
  content rather than confirmation count.
\item
  Synthetic bits are geometric appointments rather than valued outcomes
  --- a distinction forced by Axiom 3 (no value can be inherited without
  a rule; no value can be independent without dissipating polarization
  signals) and the Principle of Generative Economy.
\item
  Confirmation events resolve synthetic bit positions into binary values
  via a linear polarization-to-probability mapping \ensuremath{P(\bar{B}} = 1 \textbar{}
  \ensuremath{\Pi_{\psi})} = (1 + \ensuremath{\Pi_{\psi})/2} --- the unique mapping consistent with signal
  conservation, the maximum-entropy character of \ensuremath{U_{\mathrm{sh}},} and Generative
  Economy.
\item
  This mapping constitutes a structural correspondence to the Born rule
  for two-outcome projective measurement, paralleling the
  [[4,2,2]] structural correspondence of Paper 1 in status.
\item
  Each confirmation event is a probabilistic 3-input majority gate; the
  [[4,2,2]] stabilizer repair mechanism projects local errors
  toward stabilizer-satisfying configurations; the effective behavior in
  coherent regions is deterministic majority computation.
\item
  The substrate supports universal classical computation by the
  universality of \{MAJ, constant\}.
\item
  Programs --- configurations causally closed under the flop operator
  with K(C) \textless{} \textbar C\textbar{} --- are selected
  automatically by the stabilizer repair mechanism. Configurations that
  compute persist; configurations that do not compute dissolve into
  \ensuremath{U_{\mathrm{sh}}-\text{like}} noise.
\item
  The \ensuremath{\beta} = f(d, \ensuremath{K_{\psi})} prediction from Paper 1 acquires its computational
  meaning: \ensuremath{K_{\psi}} is the Kolmogorov complexity of the program running
  through a bit's causal history, and the GUP correction at the bit's
  position is the irreducibility of that program made geometrically
  visible.
\end{itemize}

The arc from Paper 1 to Paper 2 is a single unbroken chain: \ensuremath{U_{\mathrm{sh}}} provides
the maximum-entropy substrate; triadic confirmation pulls binary
outcomes from it; the FCC lattice carries those outcomes forward as wave
dynamics; the polarization field encodes their computational content;
the Born-rule structural correspondence governs their per-event
resolution; the majority gate is the irreducible computational
primitive; the stabilizer dynamics select programs from noise. No new
axioms. No new parameters. No new mechanisms.

The remaining derivations of Volume 1 --- learning and hierarchy (Paper
3), subsystem emergence (Paper 4), and simulation (Paper 5) --- proceed
from this computational substrate. The agent-side complement of the
framework, addressing the relationship between substrate-level
computation and observer experience, is the subject of Volume 2.

\section*{Acknowledgments}
\addcontentsline{toc}{section}{Acknowledgments}

The development of this volume was substantially aided by extended
collaboration with several large language models used as critical
sounding boards, pedagogical resources, formal-rigor checkers, and
external working memory across many sessions during 2024--2026.

\textbf{Gemini (Google DeepMind)} was the long-running primary
collaborator across the framework's conceptual development, providing
physics pedagogy, structural critique, and ongoing review of the ideas
that became Volume 1. The author's path into the technical literature,
and the working-out of the framework's core conceptual moves over a long
period preceding formalization, were carried by extended dialogue with
Gemini.

\textbf{Claude (Anthropic)} was the primary collaborator across the
formalization, derivation tightening, simulation work, citation
auditing, and editorial passes that brought the volume to its present
form, including the sublattice bifurcation analysis empirically verified
in Paper 5 Appendix E.

\textbf{ChatGPT (OpenAI)}, \textbf{DeepSeek}, \textbf{Grok (xAI)}, and
\textbf{Granite (IBM)} contributed additional structural feedback and
review at various points.

All theoretical commitments, the choice of axioms, the Principle of
Generative Economy, the sublattice bifurcation insight, and the
framework's overall architecture are the author's; the AI systems' role
was to teach, formalize, critique, verify, and document.

\section*{References}
\addcontentsline{toc}{section}{References}
\begin{reflist}
Almheiri, A., Dong, X., \& Harlow, D. (2015). Bulk locality and quantum
error correction in AdS/CFT. Journal of High Energy Physics, 2015(4),
163. doi:10.1007/JHEP04(2015)163. \url{https://arxiv.org/abs/1411.7041}

Bombelli, L., Lee, J., Meyer, D., \& Sorkin, R. D. (1987). Space-time as
a causal set. Physical Review Letters, 59(5), 521--524.
doi:10.1103/PhysRevLett.59.521.

Born, M. (1926). Zur Quantenmechanik der Stoßvorgänge. Zeitschrift für
Physik, 37(12), 863--867. doi:10.1007/BF01397477.

Goldstein, H., Poole, C., \& Safko, J. (2002). Classical Mechanics (3rd
ed.). Addison-Wesley.

Gorard, J. (2020). Some Quantum Mechanical Properties of the Wolfram
Model. Complex Systems, 29(2), 537--598.
doi:10.25088/ComplexSystems.29.2.537.

Gottesman, D. (1997). Stabilizer Codes and Quantum Error Correction.
Ph.D.~Thesis, California Institute of Technology.
\url{https://arxiv.org/abs/quant-ph/9705052}

Hoffman, D. D., Prakash, C., \& Prentner, R. (2023). Fusions of
Consciousness. Entropy, 25(1), 129. doi:10.3390/e25010129.

Hossenfelder, S. (2013). Minimal Length Scale Scenarios for Quantum
Gravity. Living Reviews in Relativity, 16(1), 2.
doi:10.12942/lrr-2013-2.

Kauffman, S. (1993). The Origins of Order. Oxford University Press.

Kolmogorov, A. N. (1965). Three approaches to the quantitative
definition of information. Problems of Information Transmission, 1(1),
1--7.

Pastawski, F., Yoshida, B., Harlow, D., \& Preskill, J. (2015).
Holographic quantum error-correcting codes: toy models for the
bulk/boundary correspondence. Journal of High Energy Physics, 2015(6),
149. doi:10.1007/JHEP06(2015)149. \url{https://arxiv.org/abs/1503.06237}

Shannon, C. E. (1948). A mathematical theory of communication. Bell
System Technical Journal, 27(3), 379--423.
doi:10.1002/j.1538-7305.1948.tb01338.x.

Surya, S. (2019). The causal set approach to quantum gravity. Living
Reviews in Relativity, 22(1), 5. doi:10.1007/s41114-019-0023-1.

Switzer, A. (2026a). The Triadic Measurement Substrate (TMS), Volume 1,
Paper 1: Emergent 3D Information Topology and the Binary Alphabet via
[[4,2,2]] Quantum Error Detection Structural Correspondence.

Takayanagi, T. (2025). Emergent Holographic Spacetime from Quantum
Information. Physical Review Letters, 134, 240001.
doi:10.1103/pg4r-fy8n.

't Hooft, G. (2016). The Cellular Automaton Interpretation of Quantum
Mechanics. Fundamental Theories of Physics 185. Springer.
doi:10.1007/978-3-319-41285-6.

Wolfram, S. (2002). A New Kind of Science. Wolfram Media.

Wolfram, S. (2020). A Class of Models with the Potential to Represent
Fundamental Physics. Complex Systems, 29(2), 107--536.
doi:10.25088/ComplexSystems.29.2.107.

Zurek, W. H. (2003). Decoherence, einselection, and the quantum origins
of the classical. Reviews of Modern Physics, 75(3), 715--775.
doi:10.1103/RevModPhys.75.715.
\end{reflist}

\clearpage


\end{document}
